% Templating thermodynamics (was the Physical model; now a subsection of the Design rules, which
% inputs this file). Brief: briefs/design.md (templating-thermodynamics part). Claims carried: the
% entries of ../claims.yaml with section model. Prose: P0-W8.3, adapted at PH-W8.2 to a mineral-neutral
% form (the nucleus is any registered mineral phase, written per formula unit); the carbonate stability
% step is replaced by the conditions K1 and K2 of the four-conditions subsection. Written to the
% derivation as audited at cp2 and cp3 (theory/derivation_body.tex; theory/thresholds.md) and the
% generalised supersaturation of contract v5. One sentence per line; each ends with a comment naming its
% ledger id. Each link is argued physically before its relation. No number appears here, and no result.
\subsection{Templating thermodynamics}\label{sec:model}

% (1) The two regimes, and which links apply in each.
The mechanism is assembled in one environment and must survive in another. % CL-R09
The template is meant to anchor, bind \ce{Ca^2+} and template the mineral while the electrode is a water-processed, alkaline slurry at room temperature (\Rs); the finished interface then sits in a non-aqueous, fluorine-free electrolyte with the collector held at the anodic potential of the cathode (\Rc). % CL-R09
The two regimes differ in solvent, reference state and mode of attack, so each relation below is used only in the regime whose reference states it assumes: adhesion, \ce{Ca^2+} binding and nucleation exist only in \Rs, coverage is judged with the \Rs\ wetting energies against a thickness bound set by electron transport in \Rc, and stability is the pair of conditions $K_2$ and $K_1$ of Section~\ref{sec:conditions} (Fig.~\ref{fig:chain}). % CL-R09
The relations hold for any mineral phase, written per formula unit with its own solubility product, molar volume and nucleus--liquid energy, and the links are necessary conditions in series: one whose failing side carries at least the credible mass of the propagated inputs ends the chain, and one that passes shows only that its step is thermodynamically possible. % CL-R09, CL-R05

% carbonate control (frozen): the artwork of this figure shows a calcium carbonate nucleus.
\begin{figure*}[t]
  \IfFileExists{\nacreRepoRoot/figures/out/fig2.pdf}{%
    \nacreFigTwoPanelLetters
    \includegraphics[width=\textwidth]{\nacreRepoRoot/figures/out/fig2.pdf}}%
    {\nacreMissing{chain artwork, pending}}
  \caption{The links of the chain the screen tests, as an illustrative construction, not a simulated structure. In every panel the brown strip is the hydroxylated aluminium oxide, with its surface oxygens as red discs. (a) Adhesion in water: interfacial water covers the bare oxide on the left; on the right, catechol-bearing units of a generic oligomer (nitrogen blue) bind the oxide through their catechol oxygens, and the water they displace leaves the contact (short arrows). (b) \ce{Ca^2+} binding: the catechol-bearing layer (beige band) carries catecholate and carboxylate sites, a few of which hold a \ce{Ca^2+} ion (green); the solution above holds the same sparse \ce{Ca^2+} and carbonate ions at every height, because binding concentrates \ce{Ca^2+} at the sites without raising the supersaturation, which stays that of the bulk. (c) A calcium carbonate nucleus (pale blue cap) at a small contact angle on the template; the template--nucleus contact (vermillion) is where the template acts on the barrier, and \ce{Ca^2+} ions bridging sites to the nucleus belong to that contact. (d) Neighbouring nuclei coalesced into a continuous film over the template, covering the whole template--film contact (vermillion). Structures are drawn from their SMILES as pictograms: bonds as lines, oxygen red, nitrogen blue, \ce{Ca^2+} green, hydrogen white, carbon at line vertices, charges not drawn. The dashed boxes with diagonals at the upper right of each panel hold the places of detail renders still to come.}
  \label{fig:chain}
\end{figure*}

% (2) Adhesion in water: the Dupre work with water competition.
In the slurry the template never meets a dry oxide: the hydroxylated \ce{Al2O3}/\ce{AlOOH} surface carries hydrogen-bonded water and the template is hydrated, so a template--oxide contact forms only by displacing that water (Fig.~\ref{fig:chain}a). % CL-R01
The reversible work per area that separates contact from the two wetted surfaces is the Dupr\'e work in water, % CL-R01
\begin{equation}\label{eq:dupre}
\Wwet=\gow+\gpw-\gop=\Wdry+2\gw-W_\mathrm{ow}-W_\mathrm{pw},
\end{equation}
where $W_\mathrm{ow}$ and $W_\mathrm{pw}$ are the works with which water binds each surface and $\Wdry=\go+\gp-\gop$ is the adhesion against vacuum. % CL-R01
A large dry adhesion is offset by the binding of water to each surface, and a hydroxylated oxide and a hydrophilic template both bind water strongly, so $\Wwet$ can be small or negative where $\Wdry$ is large; the dry work describes a state the slurry never presents and is kept only as a comparator. % CL-R01
The bare oxide has $\Wwet=0$ and fixes the zero; $\Wwet$ is the potential of mean force of a template fragment pulled from the oxide in explicit water per contact area, which already contains the displacement penalty and the entropy of the released water, and it is a descriptor of the template, not a gate. % CL-R01, CL-D03

% (3) Ca2+ binding as a descriptor.
The template can hold \ce{Ca^2+} only if its coordinating groups compete successfully with the hydration shell of the ion. % CL-R02
The transfer free energy $\DGseq$ of one \ce{Ca^2+} from the hydrated bulk to a site, built as a thermodynamic cycle that includes the deprotonation of the site at the slurry pH and the activity corrections at its ionic strength, measures that competition, and $\DGseq<0$ means that the ion accumulates at the sites (Fig.~\ref{fig:chain}b). % CL-R02
Binding raises the local \ce{Ca^2+} concentration, but it cannot raise the driving force for nucleation. % CL-R02
The interfacial region exchanges ions with the slurry, so at equilibrium every ion carries the same electrochemical potential at a site as in the bulk: binding lowers the local standard potential of the ion by exactly as much as it raises the concentration term, and a mean electrostatic potential cancels in the neutral solid. % CL-R02
The driving force per formula unit at the template is therefore the bulk one, % CL-R02
\begin{equation}\label{eq:dmu}
\Dmu=\kB T\ln\Omb,
\end{equation}
with $\Omb$ the bulk supersaturation of Eq.~\eqref{eq:stability}, built from the free-ion activities $a_i$ of the speciation at the slurry state, the ions $\nu_i$ and structural water $n$ per formula unit, and the solubility product $\Ksp$ of the framing. % CL-R02
Binding reaches the barrier only through the interfacial free energies. % CL-R02
Bound \ce{Ca^2+} is a surface excess $\GCa=\Gs\Theta-\ds\cb$ of the template--liquid interface, with $\Theta$ the Langmuir occupancy of sites of areal density $\Gs$ in a binding layer of thickness $\ds$ and $\cb$ the free-ion concentration, and the Gibbs adsorption equation turns that excess into a change of the template--liquid energy, % CL-R02
\begin{equation}\label{eq:gibbs}
\dga=\gsl-\gslz=-\Rg T\big[\Gs\ln(1+x)-\ds\cb\big],\qquad x=\frac{\ds\cb}{\Gs}\,\e^{-\DGseq/\Rg T},
\end{equation}
which is negative whenever the excess $\GCa$ is positive; $\gslz$ is the energy of the \ce{Ca}-free template. % CL-R02
Treating the adsorbing component as \ce{Ca^2+} with its compensating charge, a neutral solute, is an approximation: it leaves out the change of the diffuse-layer free energy when binding changes the surface charge, a correction that ionic strength screens and to which the model assigns no sign. % CL-R02
\ce{Ca^2+} held at the template--crystal contact, for example bridging a carboxylate to the crystal, is not part of $\dga$: it belongs to the template--nucleus energy $\gsn$, which is taken from contact data at the slurry state (Fig.~\ref{fig:cycle}). % CL-R02
$\DGseq$ is therefore a binding descriptor, reported with the occupancy, the excess and $\dga$, and not a criterion of its own. % CL-R02, CL-D03
The consequence is plain: at a fixed template--nucleus energy, binding alone raises the nucleation barrier, because the bound ions must be stripped from the area a nucleus covers, as for an adsorbed inhibitor. % CL-R10
The template helps nucleation only if the template--nucleus energy falls by more than the adsorption lowers the template--liquid energy. % CL-R10

% carbonate control (frozen): the artwork of this figure shows a calcium carbonate nucleus.
\begin{figure*}[t]
  \IfFileExists{\nacreRepoRoot/figures/out/fig3.pdf}{%
    \includegraphics[width=\textwidth]{\nacreRepoRoot/figures/out/fig3.pdf}}%
    {\nacreMissing{binding-pathway artwork, pending}}
  \caption{Where the \ce{Ca^2+} binding of the template enters the nucleation barrier, as an illustrative construction, not a simulated structure. Top left (grey box): the bulk solution at the regime activities, with a hydrated \ce{Ca^2+} ion (green disc with six water molecules, oxygen towards the ion) and a carbonate ion; every ion here carries its bulk electrochemical potential. Lower middle: the template--liquid interface at the regime \ce{Ca^2+} activity; the catechol-bearing layer (beige band) on the hydroxylated oxide (brown strip, red discs the surface oxygens) carries catecholate and carboxylate sites, two of which hold a \ce{Ca^2+} ion. Blue curved arrow, top left to lower middle: the transfer of \ce{Ca^2+} from the hydrated bulk to a site, whose free energy changes only the template--liquid energy, through the Gibbs--Langmuir term (blue line). Black arrow, lower middle to lower right: a nucleus forms on the template. Lower right: the calcium carbonate nucleus as a spherical cap (pale blue) at a small contact angle, where three interfaces meet: the template--nucleus contact under the cap (vermillion), whose \ce{Ca^2+} ions bridging sites to the crystal are its own \ce{Ca^2+} excess, counted inside its energy; the nucleus--liquid surface of the cap (bold black outline); and the template--liquid interface outside the cap (blue), still carrying its adsorbed \ce{Ca^2+}. Three short arrows at the right edge of the cap: the three interfacial energies as line tensions in Young's relation, blue outward along the template, vermillion inward along the template and black along the cap surface. Black curved arrow, top left to the cap: the formula units of the nucleus come from the bulk solution at the bulk chemical potentials, so the driving force is that of the bulk and none of it comes from the bound layer. Structures are drawn from their SMILES as pictograms, with the colours of Fig.~\ref{fig:chain}.}
  \label{fig:cycle}
\end{figure*}

% (4) Heterogeneous nucleation: Young with the fixed sign convention, single entry, absolute barrier.
A cluster of the mineral gains $\Dmu$ for every formula unit it holds and pays an interfacial free energy for every area it creates; volume grows faster than surface, so the barrier sits at a critical size. % CL-R03
On the template a nucleus replaces template--liquid interface by template--nucleus interface over its base, which lowers the barrier without changing the driving force (Fig.~\ref{fig:chain}c). % CL-R03
In the capillarity approximation for a spherical cap, % CL-R03
\begin{equation}\label{eq:het}
\begin{gathered}
\DGhet=f(\theta)\,\frac{16\pi}{3}\frac{\gnl^{3}\vm^{2}}{\Dmu^{2}},\qquad f(\theta)=\frac{(2+\cos\theta)(1-\cos\theta)^{2}}{4},\\
\cos\theta=\frac{\gslz+\dga-\gsn}{\gnl},
\end{gathered}
\end{equation}
with $\vm$ the volume per formula unit and $\Dmu$ the bulk driving force of Eq.~\eqref{eq:dmu}. % CL-R03
The sign convention is fixed: $\theta$ is measured through the nucleus, and $\cos\theta$ is the free energy gained per area by replacing template--liquid interface with template--nucleus interface, in units of the nucleus--liquid energy, so $\theta=0$ is complete wetting with no barrier and $\theta=\pi$ is no wetting and the homogeneous barrier. % CL-R03
Young's relation carries the bare $\gnl$, the same energy that sets the homogeneous barrier: water displaced from under the contact sits on the contact disc and never touches the free surface of the cap, so a displacement cost enters $\gsn$ once, as a positive correction that raises $\theta$, and not at all when $\gsn$ is already a full in-water contact energy. % CL-R03, CL-R05
Every template-dependent input enters once, in one interfacial energy, and only through $\theta$: $\Omb$ depends on bulk activities alone, while $\DGseq$, $\Gs$ and $\ds$ enter only through $\dga$, added once to the \ce{Ca}-free $\gslz$; crediting the bound ions to the driving force instead would give binding the opposite effect on the barrier, and crediting them to both would count the binding free energy twice (proof in the Supplementary Information, \nacreSIref{si:no-double-counting}{No double counting of interfacial terms}). % CL-R05
A barrier is low enough only relative to the number of attempts the system makes in the time available. % CL-R11
Nuclei form on the template at a rate per area $J=J_0\e^{-\DGhet/\kB T}$, so over the time $t_\mathrm{p}$ the template spends at the slurry states the areal density of nuclei must reach a required density, % CL-R11
\begin{equation}\label{eq:access}
\DGhet\le B_\mathrm{acc}=\kB T\ln\frac{J_0t_\mathrm{p}}{N_\mathrm{req}},\qquad N_\mathrm{req}=\frac{\ln(1/q_0)}{A_\mathrm{patch}}.
\end{equation}
The required density follows from the contact patch $A_\mathrm{patch}$ through which a cathode particle meets the collector: if nuclei land independently, the number under one patch is Poisson distributed, and allowing at most a fraction $q_0$ of the patches to stay empty gives the second relation. % CL-R11
Only $\ln J_0$ enters, so an order-of-magnitude prefactor moves $B_\mathrm{acc}$ by a few $\kB T$; the count ignores exclusion zones around growing islands, so it is an upper estimate and the criterion is necessary, not sufficient. % CL-R11

% (5) Film coverage: spreading, and why complete wetting is excluded when gamma_sl <= gamma_nl.
A low barrier says where nuclei form easily, not whether they make a continuous film, and a single gap leaves a path from the electrolyte to the collector. % CL-R12
Whether the mineral spreads over the template or beads up on it is a wetting balance, the spreading coefficient % CL-R12
\begin{equation}\label{eq:spread}
\Ssp=\gsl-\gsn-\gnl\le\gsl-\gnl,
\end{equation}
where $\Ssp\ge0$ means complete wetting and growth layer by layer and $\Ssp<0$ means growth by islands. % CL-R12
The bound holds because $\gsn\ge0$: a negative interfacial free energy between two immiscible solids would make them create contact without limit, so no sharp contact would exist. % CL-R12
Complete wetting therefore needs $\gsl\ge\gnl$, and a template that water wets at least as favourably as it wets the mineral, $\gsl\le\gnl$, cannot be wetted completely by a solid-like nucleus, whatever its contact with the nucleus; \ce{Ca^2+} binding pushes the same way by lowering $\gsl$. % CL-R12
Hydrophilic is meant here in that relative sense: a small water contact angle in air does not by itself place $\gsl$ below $\gnl$, so whether a template meets the condition is a question for the inputs, and the screen reports the prior mass with $\gsl>\gnl$. % CL-R12
On such a template a continuous film can form only when islands meet before they grow too thick (Fig.~\ref{fig:chain}d): for equal caps on a square lattice, closure below an admissible thickness $d_\mathrm{max}$ needs an island density that rises steeply with $\theta$ and with the inverse square of $d_\mathrm{max}$, which is set by electron transport through the insulating film in the cell. % CL-R12
A liquid precursor is the exception the route layer uses: its tension against the dilute liquid is small, so the bound above no longer forbids spreading (Section~\ref{sec:routes}). % CL-R12, CL-D16

% (6) Stability: carried by the conditions.
Whether a film, once formed, persists is not a templating question: dissolution in the slurry and oxidation in the cell are the conditions $K_2$ and $K_1$ of Section~\ref{sec:conditions}, and dissolution is decided by the same supersaturation $\Omb$ that drives nucleation, so a slurry state that cannot nucleate the mineral is one in which a film of it dissolves, and binding at the template changes neither. % CL-D05

% (7) What the model leaves out (stated briefly; the full list is in the Discussion).
The model leaves out four things, each taken up in the Discussion: the validity of classical nucleation theory at the size of a nucleus, with its capillarity approximation, isotropic $\gnl$ and lack of non-classical pathways such as prenucleation clusters or liquid precursors other than the coacervate of the route layer; the double-layer term of the adsorption energy; kinetics beyond the accessibility bound, including any local excess of driving force during drying; and the persistence of the layer, since anchoring, binding and nucleation are not computed in the cell. % CL-L06
